% SPDX-FileCopyrightText: 2025 IObundle
%
% SPDX-License-Identifier: MIT

Versat precalculates, for every path in the dataflow graph, how many cycles
it takes for data to travel from one unit to another, and automatically
delays faster paths to match slower ones. This guarantees that whenever
multiple paths converge on a single unit, every one of its inputs carries
valid data on the same cycle; units never need to wait for one another, and
the end user does not need to reason about these delays.

\begin{verbatim}
module ValidityExample(){
  Mem slow;
  Mem fast;
  PipelineRegister pipe; // Delays its input by one cycle.
  Reg result;
#
  slow -> pipe; // pipe delays slow by one cycle.
  sum = pipe + fast; // Versat delays 'fast' by one cycle to compensate,
                      // so both operands of the sum arrive together.
  sum -> result;
}
\end{verbatim}

This default behaviour is correct whenever a computation only concerns
values that arrive on the same cycle. When a computation must combine
values from different cycles, the designer can request an explicit delay
using the \texttt{\{N\}} operator, which shifts a unit's output stream by
$N$ cycles:

\begin{verbatim}
module DelayExample(){
  Mem mem;
#
  // The first valid value of mem{3} is the fourth value output by mem.
  sum = mem{0} + mem{1} + mem{2} + mem{3};
  sum -> out;
}
\end{verbatim}

Assuming \texttt{mem} is configured to output the natural numbers starting
at zero, the accelerator computes:

\begin{table}[H]
\centering
\begin{tabular}{|c|c|c|c|c|c|}
\hline
\rowcolor{iob-green}
\textbf{Cycle} & \texttt{mem\{0\}} & \texttt{mem\{1\}} & \texttt{mem\{2\}} & \texttt{mem\{3\}} & \textbf{sum} \\
\hline
0 & 0 & 1 & 2 & 3 & 6 \\
1 & 1 & 2 & 3 & 4 & 10 \\
2 & 2 & 3 & 4 & 5 & 14 \\
3 & 3 & 4 & 5 & 6 & 18 \\
\hline
\end{tabular}
\caption{Values computed by \texttt{DelayExample} over four cycles.}
\label{tab:data_validity}
\end{table}
